\begin{table}[htbp]
\caption{Distritos con peor acceso a capacidad resolutiva. Una raya en la columna de tiempo medio indica que el distrito no tiene ningun punto de demanda con tiempo interpretable: la red vial no lo sirve en absoluto}
\label{tab:brechas}
\begin{tabular}{lrrrrr}
\toprule
\# & Distrito & Departamento & Poblacion & Sin acceso vial (\%) & Media (min) \\
\midrule
1 & OCAÑA & AYACUCHO & 2738 & 100.0 & --- \\
2 & HUACAÑA & AYACUCHO & 715 & 100.0 & --- \\
3 & ANDOAS & LORETO & 8681 & 100.0 & --- \\
4 & BARRANCA & LORETO & 16273 & 100.0 & --- \\
5 & CAHUAPANAS & LORETO & 7859 & 100.0 & --- \\
6 & MORONA & LORETO & 31076 & 100.0 & --- \\
7 & PASTAZA & LORETO & 8067 & 100.0 & --- \\
8 & PARINARI & LORETO & 6732 & 100.0 & --- \\
9 & URARINAS & LORETO & 16535 & 100.0 & --- \\
10 & LAS AMAZONAS & LORETO & 9128 & 100.0 & --- \\
11 & TORRES CAUSANA & LORETO & 5084 & 100.0 & --- \\
12 & ROSA PANDURO & LORETO & 215 & 100.0 & --- \\
13 & TENIENTE MANUEL CLAVERO & LORETO & 4942 & 100.0 & --- \\
14 & ALTO TAPICHE & LORETO & 1932 & 100.0 & --- \\
15 & EMILIO SAN MARTIN & LORETO & 7193 & 100.0 & --- \\
16 & MAQUIA & LORETO & 7889 & 100.0 & --- \\
17 & YAGUAS & LORETO & 2359 & 97.5 & 4903.7 \\
18 & TENIENTE CESAR LOPEZ ROJAS & LORETO & 6912 & 89.1 & 3.2 \\
19 & YAVARI & LORETO & 14688 & 89.1 & 4919.5 \\
20 & TIGRE & LORETO & 8741 & 86.8 & 80.8 \\
\bottomrule
\end{tabular}
\end{table}
